\startcomponent ch01
\product fduthesis

\chapter{模板简介与配置}

\section{模板概述}

fduthesis-context 是一套基于 \ConTeXt\ (LMTX) 排版系统的
复旦大学学位论文模板，
严格遵循《复旦大学博士、硕士学位论文规范
（2024 年 10 月修订版）》。
模板支持以下论文类型和学位类别：

\startitemize
  \item 本科毕业论文（\type{type=bachelor}）
  \item 硕士学位论文（\type{type=master}）
  \item 博士学位论文（\type{type=doctor}）
  \item 学术学位（\type{degree=academic}）
  \item 专业学位（\type{degree=professional}）
\stopitemize

\section{环境要求}

使用本模板需要安装以下软件：

\startitemize[n]
  \item {\bf TeX Live 2024} 或更新版本（包含 \ConTeXt\ LMTX）。
  \item {\bf 中文字体}：Source Han Serif CN（思源宋体）、
    Source Han Sans CN（思源黑体）。
    如果系统未安装，可替换为 Noto Serif SC、SimSun 等字体，
    并在 \type{fduthesis_env.tex} 中修改字体配置。
  \item 安装或更换字体后，需运行 \type{mtxrun --script fonts --reload} 刷新字体数据库。
\stopitemize

\section{论文信息配置}

在 \type{fduthesis.tex} 主文件的开头，
通过 \type{\setvariables} 配置论文信息。
以下是一个硕士学术学位论文的配置示例：

\starttyping
\setvariables [thesis] [
  type=master,
  degree=academic,
  school-id={10246},
  student-id={22110000000},
  title={论文中文标题},
  title-en={English Title of the Thesis},
  author={作者姓名},
  department={院系名称},
  major={专业名称},
  supervisor={导师姓名\quad 职称},
  date={2025 年 6 月},
  keywords={关键词一，关键词二，关键词三},
  keywords-en={keyword1, keyword2, keyword3},
  clc={TP391},
  instructors={导师一\quad 教授\par 导师二\quad 副教授},
]
\stoptyping

\section{文件结构}

模板的文件组织如下：

\starttyping
fduthesis-context/
  fduthesis.tex         % 主文件（论文信息配置与内容组织）
  fduthesis_env.tex     % 环境定义（版式、字体、宏定义）
  references.bib        % 参考文献数据库
  Makefile              % 编译脚本
  figures/              % 图片目录
  frontmatter/          % 前置部分
    cover.tex           %   封面（封一）
    supervisors.tex     %   扉页（指导小组）
    toc.tex             %   目录
    abstract_cn.tex     %   中文摘要
    abstract_en.tex     %   英文摘要
    symbols.tex         %   符号与缩略语表
  bodymatter/           % 正文部分
    ch01.tex ~ ch05.tex %   各章内容
  backmatter/           % 后置部分
    appendix.tex        %   附录
    publications.tex    %   学术论文目录
    acknowledgements.tex%   致谢
    declarations.tex    %   独创性声明
\stoptyping

\section{编译方法}

在项目根目录执行以下命令编译论文：

\starttyping
make          % 单次编译
make twice    % 两次编译（更新目录和交叉引用）
make clean    % 清理辅助文件
make distclean% 清理所有生成文件
\stoptyping

也可以直接调用 \ConTeXt\ 编译：

\starttyping
context fduthesis.tex
\stoptyping

\section{论文内容顺序}

根据规范，论文内容按以下顺序编排：

\startitemize[n]
  \item 论文封面（封一）
  \item 扉页：指导小组成员名单
  \item 目录
  \item 中文摘要（关键词、中图分类号）
  \item 英文摘要（关键词、中图分类号）
  \item 符号与缩略语表（如有）
  \item 正文（引言、各章节）
  \item 参考文献
  \item 附录（如有）
  \item 攻读学位期间发表的学术论文
  \item 致谢
  \item 独创性声明与授权声明（封三）
\stopitemize

\stopcomponent
